\uuid{awQs}
\titre{Base/Dimension/Espace des fonctions}
\auteur{hanani ; mzari}
\organisation{exo7}
\licence{CC BY-NC-SA 4.0}
\datecreate{2020-01-20}
\niveau{L1}
\module{Algèbre}
\chapitre{Espace vectoriel}
\sousChapitre{Base}
\difficulte{3}
\theme{Base et dimension, Base, Dimension, Espace des fonctions}
\exo7qcmid{lille-2.espaces-vectoriels.050}

\contenu{
\begin{question}
\qclassification{106.04}
\qtags{motcle=Base/Dimension/Espace des fonctions}
Soit $n$ un entier $\geq 2$. On considère les fonctions réelles $f_1, f_2, \dots , f_n$, définies par :
$$f_1(x)=e^x,\quad f_2(x)= e^{2x},\quad \dots ,\quad f_n(x) = e^{nx}$$
et $E$ l'espace vectoriel engendré par ces fonctions. Quelles sont les assertions vraies ?
\begin{answers}
\bad{$E$ est un espace vectoriel de dimension $n-2$}.
\bad{$E$ est un espace vectoriel de dimension $n-1$}.
\good{$E$ est un espace vectoriel de dimension $n$}.
\bad{$E$ est un espace vectoriel de dimension infinie}.
\end{answers}
\begin{explanations} Soit $\lambda_1, \dots, \lambda_n$ des réels tels que $\lambda_1e^x+ \dots+ \lambda_ne^{2x}=0,$
pour tout réel $x$. En divisant par $e^x$ et en faisant tendre $x$ vers $-\infty$, on obtient $\lambda_1=0$. Puis, en divisant par $e^{2x}$ et en faisant tendre $x$ vers $-\infty$, on obtient $\lambda_2=0$. En appliquant ce raisonnement $n$ fois, on démontre que tous les $\lambda_i$ sont nuls. Par conséquent,
$\{f_1, f_2,  \dots , f_n\}$ est une base de $E$ et donc $\dim E =n$.
\end{explanations}
\end{question}
}
